\section{Broader impacts}
\label{appendix:impacts}
This work introduces CADS, a compute-aware data selection algorithm that dynamically adapts the training data budget according to available computational resources. By optimizing data efficiency relative to compute constraints, CADS significantly reduces the computational overhead of training deep learning models without sacrificing performance. This advancement enables more accessible and environmentally sustainable machine learning research and deployment, particularly in scenarios constrained by limited hardware resources.
The method’s potential to lower energy consumption during model training aligns with broader community efforts towards greener AI, contributing to reductions in carbon footprint and operational costs. Additionally, CADS facilitates democratization of deep learning by empowering researchers and practitioners with modest resources to train competitive models.
However, as with any automation that accelerates model training, there exist potential risks, including the possibility of speeding up the development of models with harmful biases, privacy vulnerabilities, or malicious intent if not applied responsibly. We stress the importance of ethical use and compliance with community norms and regulatory standards when deploying such techniques.
Overall, CADS represents a step forward in efficient and responsible machine learning practices, promoting sustainability and equitable access within the deep learning community.